
\documentclass[11pt,a4paper]{article}

\usepackage[margin=0.8in]{geometry}
\usepackage[T1]{fontenc}
\usepackage{lmodern}
\usepackage{xcolor}
\usepackage{hyperref}
\usepackage{enumitem}
\usepackage{titlesec}
\usepackage{tabularx}
\usepackage{fancyhdr}
\usepackage{tcolorbox}
\usepackage{amssymb}

\definecolor{darkblue}{RGB}{30,65,110}
\definecolor{lightblue}{RGB}{235,243,252}
\definecolor{lightgray}{RGB}{245,245,245}
\definecolor{darkgreen}{RGB}{35,105,70}

\hypersetup{
    colorlinks=true,
    linkcolor=darkblue,
    urlcolor=darkblue
}

\titleformat{\section}
    {\Large\bfseries\color{darkblue}}
    {\thesection}{0.6em}{}

\titleformat{\subsection}
    {\large\bfseries\color{darkblue}}
    {\thesubsection}{0.5em}{}

\titleformat{\subsubsection}
    {\normalsize\bfseries\color{darkgreen}}
    {\thesubsubsection}{0.5em}{}

\pagestyle{fancy}
\fancyhf{}
\lhead{Unit 1 -- Clean Coding}
\rhead{Course Tracker}
\cfoot{\thepage}

\newtcolorbox{conceptbox}{
    colback=lightblue,
    colframe=darkblue,
    boxrule=0.7pt,
    arc=2mm,
    left=2mm,right=2mm,top=1.5mm,bottom=1.5mm
}

\newtcolorbox{exambox}{
    colback=lightgray,
    colframe=black!40,
    boxrule=0.5pt,
    arc=2mm,
    left=2mm,right=2mm,top=1.5mm,bottom=1.5mm
}

\newcommand{\checkitem}{$\square$}
\newcommand{\code}[1]{\texttt{#1}}

\title{
    \vspace{-1cm}
    \textbf{Unit 1 -- Clean Coding}\\
    \large Complete Topic \& Subtopic Tracker
}
\author{Course Notes}
\date{}

\begin{document}

\maketitle

\begin{conceptbox}
\textbf{Syllabus scope:} This document covers the ``Clean coding mindset'' portion of Unit 1:
readable code, modular code, maintainable code, naming conventions, avoiding code smells,
and scalable code structure.

\vspace{2mm}

\textbf{How to use this document:} Treat each subsection as an individual study item.
Mark the checkbox only after you understand the topic and can explain it in your own words.
\end{conceptbox}

\tableofcontents
\newpage

% =========================================================
\section{Clean Coding Mindset}

\subsection{Study Tracker}

\begin{itemize}[label=\checkitem]
    \item Understand what clean coding means
    \item Understand why code quality matters
    \item Understand readability
    \item Understand modularity
    \item Understand maintainability
    \item Understand the relationship between readability, modularity, and maintainability
\end{itemize}

\subsection{What is Clean Code?}

Clean code is code that is structured and written in a way that makes its purpose
and behavior relatively easy for developers to understand and modify.

The syllabus describes the clean coding mindset through these qualities:

\begin{itemize}
    \item Readable code
    \item Modular code
    \item Maintainable code
    \item Clear naming
    \item Avoidance of code smells
    \item Scalable code structure
\end{itemize}

\subsection{Why Clean Coding Matters}

Poorly structured code can become increasingly difficult to understand and modify
as a project grows.

Clean coding aims to make the codebase easier to:

\begin{itemize}
    \item Read
    \item Understand
    \item Debug
    \item Modify
    \item Extend
    \item Maintain
\end{itemize}

\begin{exambox}
\textbf{Exam focus:} Be able to explain why clean code is important rather than
defining it only as ``good-looking code.''
\end{exambox}

\subsection{Clean Code vs Working Code}

A program can produce the correct output while still having poor structure.

\begin{center}
\textbf{Correct behavior} $\neq$ \textbf{Good design}
\end{center}

For example, code may work but still contain:

\begin{itemize}
    \item unclear variable names,
    \item duplicated logic,
    \item very large functions,
    \item tightly connected components,
    \item difficult-to-change structure.
\end{itemize}

The clean coding mindset considers both correctness and the quality of the
code structure.

% =========================================================
\section{Readable Code}

\subsection{Study Tracker}

\begin{itemize}[label=\checkitem]
    \item Understand readability
    \item Use meaningful names
    \item Keep logic understandable
    \item Avoid unnecessary complexity
    \item Use consistent formatting
    \item Understand why readability helps maintenance
\end{itemize}

\subsection{Meaningful Names}

Names should communicate what a variable, function, class, or object represents.

Prefer:

\begin{verbatim}
studentCount
calculateSalary()
PaymentService
\end{verbatim}

over vague names such as:

\begin{verbatim}
x
doStuff()
PS
\end{verbatim}

A good name reduces the amount of explanation needed to understand the code.

\subsection{Readable Functions}

A function should be understandable without requiring the reader to mentally
trace many unrelated operations.

When a function becomes too large or performs unrelated tasks, consider
splitting it into smaller functions with clear responsibilities.

\subsection{Avoid Unnecessary Complexity}

Code should solve the problem without introducing unnecessary layers or
complicated logic.

Questions to ask:

\begin{itemize}
    \item Can another programmer understand this quickly?
    \item Is the logic more complicated than the requirement?
    \item Are there unnecessary conditions or nested blocks?
    \item Does the function contain unrelated responsibilities?
\end{itemize}

\subsection{Consistent Formatting}

Consistent formatting makes code easier to scan.

Examples include:

\begin{itemize}
    \item consistent indentation,
    \item consistent naming style,
    \item consistent placement of braces,
    \item consistent spacing.
\end{itemize}

\begin{exambox}
\textbf{Key idea:} Readability is primarily about making the intent and structure
of the code easy for another developer to understand.
\end{exambox}

% =========================================================
\section{Modular Code}

\subsection{Study Tracker}

\begin{itemize}[label=\checkitem]
    \item Understand modularity
    \item Understand modules/components
    \item Separate responsibilities
    \item Understand interfaces between modules
    \item Understand benefits of modular design
\end{itemize}

\subsection{What is Modularity?}

Modularity means dividing a larger system into smaller components or modules,
where each module has a clear purpose.

Instead of:

\begin{center}
\textbf{One giant program}
\end{center}

use:

\[
\text{System}
\rightarrow
\begin{cases}
\text{User Module}\\
\text{Payment Module}\\
\text{Notification Module}\\
\text{Data Module}
\end{cases}
\]

The exact modules depend on the requirements of the system.

\subsection{Separation of Responsibilities}

A module should focus on a clearly defined part of the system.

For example, in a simple application:

\begin{verbatim}
PaymentService
NotificationService
UserService
\end{verbatim}

would have more clearly separated responsibilities than one class handling
payment, email, users, and database operations simultaneously.

\subsection{Benefits of Modularity}

Modularity can make a system:

\begin{itemize}
    \item easier to understand,
    \item easier to test,
    \item easier to modify,
    \item easier to debug,
    \item easier to extend,
    \item easier for multiple developers to work on.
\end{itemize}

\subsection{Module Boundaries}

A good module should have a clear purpose and a clear way for other modules
to interact with it.

Think of a module as a component with:

\[
\boxed{\text{Input} \rightarrow \text{Module} \rightarrow \text{Output}}
\]

Other parts of the program should not need to understand every internal detail
of that module.

% =========================================================
\section{Maintainable Code}

\subsection{Study Tracker}

\begin{itemize}[label=\checkitem]
    \item Understand maintainability
    \item Understand why requirements change
    \item Understand ease of modification
    \item Understand duplication and maintenance cost
    \item Understand how clean structure supports maintenance
\end{itemize}

\subsection{What is Maintainability?}

Maintainability describes how easily software can be understood, corrected,
modified, or extended after it has been written.

A maintainable system makes future changes easier.

\subsection{Why Maintainability Matters}

Software rarely remains unchanged.

Requirements may change because of:

\begin{itemize}
    \item new features,
    \item changed business rules,
    \item bug fixes,
    \item new integrations,
    \item changes in user requirements.
\end{itemize}

If the code is tightly coupled or poorly organized, a small change can require
changes throughout the system.

\subsection{Characteristics of Maintainable Code}

Maintainable code generally has:

\begin{itemize}
    \item meaningful names,
    \item clear responsibilities,
    \item low unnecessary duplication,
    \item understandable structure,
    \item manageable modules,
    \item predictable behavior.
\end{itemize}

\subsection{Changeability}

A useful question when reviewing a design is:

\begin{center}
\textit{``If this requirement changes, how many unrelated parts of the code
will I have to modify?''}
\end{center}

A design that localizes changes is generally easier to maintain.

% =========================================================
\section{Naming Conventions}

\subsection{Study Tracker}

\begin{itemize}[label=\checkitem]
    \item Understand naming conventions
    \item Variable naming
    \item Function/method naming
    \item Class naming
    \item Constant naming
    \item Avoid vague names
    \item Maintain naming consistency
\end{itemize}

\subsection{Variable Names}

Variable names should communicate their purpose.

Good examples:

\begin{verbatim}
totalAmount
studentCount
currentBalance
\end{verbatim}

Poor examples:

\begin{verbatim}
x
a
temp1
data
\end{verbatim}

A generic name is not always wrong, but it should have a clear reason for
being used.

\subsection{Function Names}

Function names should communicate an action or operation.

Examples:

\begin{verbatim}
calculateTotal()
findStudent()
sendNotification()
validatePayment()
\end{verbatim}

The name should help the reader understand what the function does.

\subsection{Class Names}

Class names generally represent an object, concept, or responsibility.

Examples:

\begin{verbatim}
Student
Payment
Order
PaymentService
ParkingSpot
\end{verbatim}

\subsection{Constants}

Constants should be distinguishable from ordinary variables according to the
naming convention used by the language/project.

For example:

\begin{verbatim}
MAX_RETRIES
DEFAULT_TIMEOUT
\end{verbatim}

\subsection{Consistency}

A project should use a consistent naming style.

For example, do not mix:

\begin{verbatim}
student_name
studentName
StudentName
\end{verbatim}

without a reason.

\begin{exambox}
\textbf{Exam focus:} Naming conventions are not merely cosmetic. Clear names
reduce ambiguity and make code easier to understand and maintain.
\end{exambox}

% =========================================================
\section{Avoiding Code Smells}

\subsection{Study Tracker}

\begin{itemize}[label=\checkitem]
    \item Understand what a code smell is
    \item Recognize long methods
    \item Recognize large/god classes
    \item Recognize duplicated code
    \item Recognize long parameter lists
    \item Recognize deep nesting
    \item Recognize magic numbers
    \item Understand refactoring as a response
\end{itemize}

\subsection{What is a Code Smell?}

A code smell is a sign or symptom in code that may indicate a deeper design
or maintainability problem.

A smell is not necessarily a bug.

\begin{center}
\textbf{Code smell} $\neq$ \textbf{runtime error}
\end{center}

A program can work correctly and still contain code smells.

\subsection{Long Method}

A method becomes difficult to understand because it contains too much logic.

Possible response:

\begin{itemize}
    \item identify separate responsibilities,
    \item extract smaller methods,
    \item give each method a clear purpose.
\end{itemize}

\subsection{Large or God Class}

A class becomes responsible for too many unrelated parts of the system.

Example:

\begin{verbatim}
ApplicationManager
 - users
 - payments
 - emails
 - reports
 - database
 - authentication
 - logging
\end{verbatim}

A possible refactoring is to divide these responsibilities into focused
classes/modules.

\subsection{Duplicated Code}

The same or very similar logic appears in multiple places.

Problem:

\[
\text{Change required}
\rightarrow
\text{many copies must be updated}
\]

This increases maintenance effort and can lead to inconsistent behavior.

\subsection{Long Parameter List}

A function requiring many parameters can become difficult to understand and
call correctly.

For example:

\begin{verbatim}
createUser(name, age, city, country, phone, email, ...)
\end{verbatim}

A possible design improvement is to group related information into an
appropriate object when that makes sense.

\subsection{Deep Nesting}

Many nested conditions or loops can make control flow difficult to follow.

Example structure:

\begin{verbatim}
if
  if
    for
      if
        ...
\end{verbatim}

Possible improvements include simplifying conditions, extracting functions,
or using guard clauses where appropriate.

\subsection{Magic Numbers}

A literal number appears in code without explaining its meaning.

Example:

\begin{verbatim}
if (age > 18)
\end{verbatim}

Depending on the context, a named constant can communicate intent more clearly:

\begin{verbatim}
MINIMUM_ADULT_AGE
\end{verbatim}

\subsection{Refactoring}

Refactoring means improving the internal structure of code without changing
its intended external behavior.

Typical goals:

\begin{itemize}
    \item remove duplication,
    \item simplify complex code,
    \item improve names,
    \item separate responsibilities,
    \item improve modularity.
\end{itemize}

\begin{exambox}
\textbf{Important distinction:}\\
Fixing a bug changes incorrect behavior.\\
Refactoring changes the internal structure while preserving intended behavior.
\end{exambox}

% =========================================================
\section{Scalable Code Structure}

\subsection{Study Tracker}

\begin{itemize}[label=\checkitem]
    \item Understand scalable structure
    \item Understand separation of responsibilities
    \item Understand loose coupling
    \item Understand cohesive modules
    \item Understand extensibility
    \item Understand why structure matters as projects grow
\end{itemize}

\subsection{What Does Scalable Code Structure Mean?}

In the context of the syllabus, scalable code structure means organizing code
so that the system can grow in functionality and complexity without becoming
unreasonably difficult to understand or modify.

\subsection{Separation of Responsibilities}

Separate different responsibilities instead of placing all functionality into
one component.

For example:

\[
\text{Application}
\rightarrow
\begin{cases}
\text{User Management}\\
\text{Payment}\\
\text{Notification}\\
\text{Storage}
\end{cases}
\]

Each component can then evolve more independently.

\subsection{Cohesion}

Cohesion describes how closely related the responsibilities inside a module
are.

A highly cohesive module focuses on one related area.

\begin{center}
\textbf{High cohesion} $\rightarrow$ related responsibilities stay together.
\end{center}

\subsection{Coupling}

Coupling describes the degree of dependency between modules.

A design with unnecessary strong dependencies can become difficult to change.

\begin{center}
\textbf{Good general goal: High cohesion + controlled/low coupling}
\end{center}

\subsection{Extensibility}

A scalable structure should make it possible to add functionality without
rewriting unrelated parts of the system.

For example, if a payment system supports multiple payment methods, the design
should ideally allow a new method to be introduced without modifying every
existing payment-related component.

This idea connects directly to the SOLID principles in Unit 2, especially
the Open/Closed Principle.

% =========================================================
\section{Clean Coding -- Connections Between Topics}

The topics in this section should not be memorized independently.

\subsection{How the Concepts Connect}

\[
\boxed{
\text{Good Names}
+
\text{Readable Functions}
+
\text{Modular Design}
+
\text{Clear Responsibilities}
}
\]

\[
\Downarrow
\]

\[
\boxed{\text{More Maintainable Code}}
\]

Code smells often indicate that one or more of these qualities is missing.

For example:

\[
\text{God Class}
\rightarrow
\text{poor separation of responsibilities}
\rightarrow
\text{low cohesion}
\rightarrow
\text{harder maintenance}
\]

\subsection{Clean Coding and OOP}

Clean coding provides the mindset for writing understandable code.

OOP then provides concepts such as:

\begin{itemize}
    \item encapsulation,
    \item abstraction,
    \item inheritance,
    \item composition,
    \item polymorphism.
\end{itemize}

These concepts can be used to build modular and maintainable systems.

\subsection{Clean Coding and SOLID}

Unit 1 prepares the foundation for Unit 2.

For example:

\begin{verbatim}
Large class
    |
    +--> too many responsibilities
    |
    +--> SRP violation
\end{verbatim}

Similarly:

\begin{verbatim}
Hard-coded concrete dependency
    |
    +--> difficult to replace/test
    |
    +--> DIP becomes relevant
\end{verbatim}

Thus, SOLID can be viewed as a more formal set of design principles that helps
address common maintainability and design problems.

% =========================================================
\section{Quick Revision Sheet}

\subsection{Definitions to Know}

\begin{itemize}
    \item \textbf{Clean Code:} Code designed to be understandable, maintainable,
    and well structured.
    \item \textbf{Readable Code:} Code whose intent and structure are easy to understand.
    \item \textbf{Modularity:} Dividing a system into focused components/modules.
    \item \textbf{Maintainability:} Ease with which software can be understood,
    corrected, modified, or extended.
    \item \textbf{Code Smell:} A symptom suggesting a possible design or
    maintainability problem.
    \item \textbf{Refactoring:} Improving internal code structure while preserving
    intended behavior.
    \item \textbf{Cohesion:} How closely related the responsibilities within a
    module are.
    \item \textbf{Coupling:} The degree of dependency between modules.
\end{itemize}

\subsection{Final Checklist}

\begin{itemize}[label=\checkitem]
    \item I can explain the clean coding mindset.
    \item I can explain readable code.
    \item I can identify poor naming.
    \item I can explain modular code.
    \item I can explain maintainability.
    \item I can identify common code smells.
    \item I can explain refactoring.
    \item I can explain scalable code structure.
    \item I understand cohesion.
    \item I understand coupling.
    \item I can explain how these concepts connect to SOLID.
\end{itemize}

\begin{conceptbox}
\textbf{One-line mental model:}

\[
\boxed{
\text{Clear names}
\rightarrow
\text{clear responsibilities}
\rightarrow
\text{modular structure}
\rightarrow
\text{maintainable code}
}
\]

When you encounter a piece of code, ask:
\textit{``Can another developer understand this, change it, and extend it
without having to understand the entire system?''}
\end{conceptbox}

\end{document}
